\chapter{Zusammenfassung}
\label{ch:abstract_de}

Systeme mit miteinander verbundenen Komponenten werden häufig als Netzwerke modelliert, sei es in den Naturwissenschaften wie Biologie und Physik oder in den Sozialwissenschaften. Von besonderem Interesse sind dabei Transportprozesse in diesen kontrollierten Systemen oder Netzwerken – die häufig mithilfe von Techniken zur Flussoptimierung untersucht werden. Diese Arbeit konzentriert sich auf den Transport in Netzwerken mit Reaktionen, die alternativ als \emph{Reaktionsnetzwerke} bezeichnet werden, wobei die Strukturen der miteinander wechselwirkenden Moleküle bekannt sind. Diese Struktur beeinflusst die im Netzwerk identifizierten Pfade. Reaktionsnetzwerke lassen sich mathematisch gut als Hypergraphen darstellen, was die Suche nach Pfaden mit rechnerischen Mitteln erleichtert.

Die Größe, die komplexe Vernetzung und die kombinatorische Komplexität von Reaktionsnetzwerken machen es schwierig, die darin ablaufenden Transportprozesse manuell zu analysieren. Insbesondere die Suche nach möglichen Reaktionswegen zu einem bestimmten Zielmolekül ausgehend von einer Reihe verfügbarer Ausgangsmoleküle, die gleichzeitig den Naturgesetzen entsprechen, ist eine häufig anzutreffende Fragestellung.
Diese Suche nach Reaktionswegen, die die Bildung des Zielmoleküls optimieren, ist ein zentrales Problem in vielen Anwendungsbereichen, darunter auch in biochemischen Reaktoren.
Diese Arbeit integriert etablierte Berechnungswerkzeuge aus der Chemie und physikalisch fundierte Modellierung mit bestehenden algorithmischen Methoden zur Ermittlung von Pfaden in Reaktionsnetzwerken. Sie erweitert die verfügbare ausführbare Graphentransformationssprache weiter, um die Erforschung von Reaktionsnetzwerken näher an die Realität heranzuführen.
Insbesonder präsentiert sie eine Formulierung des Suchproblems für Reaktionswege mittels gemischt-ganzzahliger linearer Programmierung (mixed ILP), bei der die chemischen Potenziale und Konzentrationen der einzelnen Moleküle integriert sind, da die Struktur bekannt ist. Dies ermöglicht es, die Suche so einzuschränken, dass nur Reaktionswege zurückgegeben werden, die ausschließlich thermodynamisch günstige Reaktionen enthalten. Werden zudem mehrere mögliche Reaktionswege gefunden, können diese anhand einer auf Thermodynamik basierenden Zielfunktion in eine Rangfolge gebracht werden. 

Die auf ganzzahliger linearer Programmierung basierende Berechnungsmethodik zur Suche nach Synthesewegen zu Zielmolekülen in Reaktionsnetzwerken wird durch die Annotation der Netzwerke mit Informationen zur Kinetik der einzelnen Reaktionen weiter erweitert. 
Die Kinetik der Reaktionen wurde anhand eines ausgewählten Modells der computergestützten Chemie aus der Struktur der Moleküle abgeleitet.
Häufig entsprechen mehrere Wege den vorgegebenen Suchkriterien. Um diese zu bewerten, wird eine Zielfunktion entwickelt, die auf physikalischen Argumenten basiert und die Wahrscheinlichkeit des jeweiligen Weges maximiert. Dies führt zu einer automatisierten Pipeline, die eine flexible und kinetisch fundierte Untersuchung von Reaktionswegen in großen Reaktionsnetzwerken mit rechnerischen Mitteln ermöglicht – selbst für Netzwerke ohne kinetische Annotation, wie sie beispielsweise durch generative Ansätze zur Erweiterung molekularer Räume erzeugt werden. 

Häufig sind die Reaktionsnetzwerke, in denen nach Reaktionswegen gesucht werden soll, zu groß, um als explizite molekulare Netzwerke generiert und gespeichert zu werden. Für solche großen Reaktionsnetzwerke, die implizit durch eine Graphgrammatik definiert sind, wird eine stochastische explorative Methode vorgestellt. Diese explorative Strategie  ahmt die stochastische chemische Kinetik nach, indem sie eine kollisionsbasierte Paarauswahl mit der Regelanwendung einer Graphgrammatik kombiniert. In jedem Schritt wählt der Algorithmus zunächst Moleküle aus, die kollidieren sollen, wählt dann eine Reaktionsregel aus und schließlich eine konkrete Reaktionsinstanz aus den möglichen Passungen der Reaktionsregel auf die beiden Moleküle. Diese „Collision-First“-Strategie vermeidet eine erschöpfende Aufzählung aller derzeit möglichen Reaktionen und ermöglicht die Exploration großer Reaktionsnetzwerke unter offenen oder geschlossenen Systembedingungen. Das erkundete Teilnetzwerk kann dann als repräsentativ für den gesamten molekularen Raum betrachtet und mit den oben genannten Methoden analysiert werden. 

Die neuen Methoden werden durch Fallstudien ergänzt, die ihre jeweiligen Stärken und Schwächen veranschaulichen. 
Die Integration der Thermodynamik wurde auf ein Reaktionsnetzwerk angewendet, das die Chemie von Blausäure und Formamid beschreibt. Alternative Reaktionswege zu den derzeit in der Literatur angenommenen werden untersucht und aufgezählt, darunter auch solche, die gemäß der gewählten Zielfunktion besser sind. 
Ein chemisches Reaktionsnetzwerk, das aus 2-Hydroxyethannitril (Glykolnitril), Wasser und Ammoniak generiert wurde, wurde mit Energiebarrieren annotiert und mit Hilfe der neuen Methoden Reaktionswegen zu Glycin und 2-Hydroxyessigsäure (Glykolsäure) gesucht.
Die neuen Methoden zur stochastischen Erkundung großer Molekülräume werden anhand der Formose-Chemie als Fallstudie veranschaulicht. Dabei werden sowohl das chemische Verhalten in dem durch die Erkundung erreichten Teilnetzwerk als auch die rechnerischen Auswirkungen des Caching auf den Explorationsprozess analysiert.
